\documentclass[journal]{IEEEtran}
\usepackage{graphicx}
\usepackage{amsmath}
\usepackage{listings}
\usepackage{xcolor}
\usepackage{hyperref}
\usepackage{booktabs}

\lstset{
  language=Python,
  basicstyle=\ttfamily\small,
  keywordstyle=\color{blue},
  commentstyle=\color{gray},
  stringstyle=\color{orange},
  breaklines=true,
  frame=single,
  numbers=left,
  numberstyle=\tiny\color{gray}
}

\title{Noise Filtering for Natural Image Restoration:\\Scenery Image}
\author{Neal Kauntia\\ARK Fresher Selection Task 2.2}

\begin{document}
\maketitle

\begin{abstract}
This document describes the approach taken to restore a heavily corrupted colour scenery 
photograph (\texttt{noisy.jpg}) degraded by dense multicolour salt-and-pepper noise. 
Unlike a binary line drawing, a natural colour image requires preservation of smooth 
colour gradients, soft textures, and structural boundaries simultaneously. A pipeline 
combining Gaussian pre-smoothing, Non-Local Means (NLM) denoising, and morphological 
opening was developed. Extensive parameter tuning was performed at each stage. The 
final output (\texttt{scenary\_final.png}) demonstrates strong noise suppression with 
good retention of the sky gradient, horizon line, and landscape textures.
\end{abstract}

\section{Introduction}
Natural image denoising differs fundamentally from line drawing denoising. A scenery 
photograph contains continuous-tone regions with smooth gradients, fine texture 
patterns, and soft edges — all of which must be preserved under denoising. The noise 
in \texttt{noisy.jpg} is multicolour (RGB channel-independent salt-and-pepper noise), 
which means each channel must be treated, and cross-channel correlations must be 
respected to avoid colour fringing artefacts.

This document covers the problem analysis, all intermediate approaches explored, 
parameter tuning decisions, and the rationale for the final chosen method.

\section{Problem Statement}
The input image \texttt{noisy.jpg} shows a landscape scene — sky gradient transitioning 
to a hilly silhouette — corrupted by very dense multicolour noise. Key characteristics:

\begin{itemize}
    \item \textbf{Noise type:} Dense salt-and-pepper noise on all three RGB channels 
    independently. Individual pixels display random saturated colours far from their 
    true value.
    \item \textbf{Signal type:} A colour photograph with smooth sky gradients 
    (cyan-to-blue transition), a teal horizontal band, and dark hill silhouettes.
    \item \textbf{Core challenge:} Noise density is so high that local filters 
    operating on small windows may still have a majority of noisy pixels in the 
    window, causing them to fail. A larger search window or non-local approach 
    is required.
    \item \textbf{Constraints:} Colour fidelity must be preserved — hue shifts 
    from channel-independent processing are unacceptable.
\end{itemize}

\section{Related Work}
\textbf{Gaussian Blur} is a linear smoothing filter. It handles 
Gaussian noise well but spreads colour noise energy across neighbouring pixels 
rather than removing it, resulting in a blurry, colour-smeared output.

\textbf{Median Blur} is highly effective for monochrome salt-and-pepper 
noise but its channel-independent application to colour images can cause hue artefacts. 
\texttt{cv.medianBlur} in OpenCV operates per-channel in colour mode, which can 
introduce colour shifts at edges.

\textbf{Non-Local Means (NLM)} searches a large region of the 
image for patches similar to the patch around the target pixel, and averages them 
weighted by patch similarity. It is particularly well-suited for high-density noise 
because its large search window finds uncorrupted similar patches even when local 
neighbours are mostly noisy. \texttt{cv.fastNlMeansDenoisingColored} applies NLM 
jointly in the LAB colour space, preserving colour consistency.

\textbf{Morphological Opening} after NLM removes any residual 
isolated bright noise clusters not fully averaged out by NLM.

\section{Initial Attempts}

\subsection{Attempt 1: Gaussian Blur Alone}
\begin{lstlisting}
result = cv.GaussianBlur(img, (k, k), sigma)
\end{lstlisting}

\textbf{Kernel and sigma tuning:}

\begin{table}[h]
\centering
\caption{Gaussian Blur Parameter Tuning}
\begin{tabular}{@{}ccll@{}}
\toprule
Kernel & $\sigma$ & Noise Reduction & Side Effects \\
\midrule
$3\times3$ & 0.5 & Very low & Colour fringing persists \\
$5\times5$ & 1.0 & Moderate & Scene blurred, edges soft \\
$7\times7$ & 1.5 & Better & Horizon line lost \\
$11\times11$ & 2.5 & High & Image unrecognisably blurry \\
\bottomrule
\end{tabular}
\end{table}

Gaussian blur distributes noise energy into neighbors rather than removing it — 
at small kernels the noise remains; at large kernels the image is destroyed. 
It was ruled out as the primary denoising method.

\subsection{Attempt 2: Median Blur with Various Kernel Sizes}
\begin{lstlisting}
result = cv.medianBlur(img, ksize)  # applied to colour image
\end{lstlisting}

\textbf{Kernel size tuning:}

\begin{table}[h]
\centering
\caption{Median Blur Kernel Size Comparison on Colour Image}
\begin{tabular}{@{}cll@{}}
\toprule
Kernel & Noise Removed & Side Effects \\
\midrule
$3\times3$ & Low — dense noise survives & Minimal damage \\
$5\times5$ & Moderate & Slight colour bleed at edges \\
$7\times7$ & Good & Noticeable texture smearing \\
$9\times9$ & Strong & Sky gradient steps/bands appear \\
$11\times11$ & Very strong & Colour posterisation, gradient lost \\
\bottomrule
\end{tabular}
\end{table}

Median filtering on its own was insufficient at the noise density present in this image. 
At $3\times3$ and $5\times5$, the noise density is so high that a majority of pixels 
in the window are still noisy, so the median remains a noisy value. Larger kernels 
introduced gradient banding and texture destruction. This motivated moving to NLM.

\subsection{Attempt 3: Non-Local Means Alone — Parameter Tuning}
\begin{lstlisting}
result = cv.fastNlMeansDenoisingColored(
    img, None,
    h=h_val, hColor=hc_val,
    templateWindowSize=tw,
    searchWindowSize=sw
)
\end{lstlisting}

\textbf{Parameter definitions:}
\begin{itemize}
    \item \texttt{h}: Filter strength for luminance. Higher = more denoising but 
    more detail loss.
    \item \texttt{hColor}: Filter strength for colour channels. Should typically 
    equal or be slightly less than \texttt{h}.
    \item \texttt{templateWindowSize}: Size of the patch used for similarity 
    comparison. Larger = more context but slower. Must be odd.
    \item \texttt{searchWindowSize}: Size of the region searched for similar 
    patches. Larger = better denoising but much slower. Must be odd.
\end{itemize}

\begin{table}[h]
\centering
\caption{NLM Parameter Tuning Results}
\begin{tabular}{@{}cccll@{}}
\toprule
h & hColor & templateW & searchW & Observation \\
\midrule
5 & 5 & 7 & 15 & Insufficient — noise remains \\
10 & 10 & 7 & 21 & Good balance — chosen base \\
15 & 10 & 7 & 21 & Over-smoothed — sky texture lost \\
10 & 10 & 5 & 21 & Slightly weaker colour recovery \\
10 & 10 & 7 & 35 & Better but very slow \\
20 & 15 & 9 & 21 & Strong noise removal, gradient slightly flat \\
\bottomrule
\end{tabular}
\end{table}

\texttt{h=10, hColor=10, templateWindowSize=7, searchWindowSize=21} provided the 
best balance: noise substantially removed, sky gradient intact, and hill silhouette 
edges preserved. Larger \texttt{h} values caused the sky gradient to flatten 
perceptibly.

\subsection{Attempt 4: Gaussian Pre-smoothing before NLM}
The observation that NLM still left some residual high-frequency colour noise 
at \texttt{h=10} led to adding a light Gaussian pre-smoothing step.

\begin{lstlisting}
gaus = cv.GaussianBlur(img, (k, k), sigma)
result = cv.fastNlMeansDenoisingColored(gaus, ...)
\end{lstlisting}

\textbf{Kernel tuning for pre-smoothing Gaussian:}

\begin{table}[h]
\centering
\caption{Gaussian Pre-smoothing Kernel Tuning}
\begin{tabular}{@{}ccl@{}}
\toprule
Kernel & $\sigma$ & Effect on NLM Input \\
\midrule
$3\times3$ & 0.5 & Subtle — slightly improves NLM patch similarity \\
$5\times5$ & 1.0 & Noticeable improvement — residual noise clusters reduced \\
$7\times7$ & 1.5 & Pre-smoothing too strong — edges enter NLM already blurred \\
\bottomrule
\end{tabular}
\end{table}

A $(3\times3)$ Gaussian with $\sigma=0.5$ was chosen: it reduces the amplitude of 
extreme noise pixels just enough to improve NLM's patch-similarity computation 
without pre-blurring the actual scene features that NLM needs to find good matches for.

\subsection{Attempt 5: Morphological Opening Kernel Tuning}
After NLM, occasional small bright clusters remain. Morphological opening was added 
as a final clean-up step.

\begin{lstlisting}
kernel = cv.getStructuringElement(cv.MORPH_ELLIPSE, (k, k))
morph = cv.morphologyEx(result, cv.MORPH_OPEN, kernel)
\end{lstlisting}

\textbf{Kernel shape and size tuning:}

\begin{table}[h]
\centering
\caption{Morphological Opening Kernel Tuning}
\begin{tabular}{@{}ccll@{}}
\toprule
Shape & Size & Noise Removed & Scene Damage \\
\midrule
RECT & $2\times2$ & Small isolated clusters & Minimal \\
ELLIPSE & $2\times2$ & Same as RECT at this size & Minimal \\
RECT & $3\times3$ & Slightly more & Faint texture erosion \\
ELLIPSE & $3\times3$ & Slightly more, rotationally symmetric & Minimal \\
ELLIPSE & $4\times4$ & Aggressive & Colour regions slightly eroded \\
\bottomrule
\end{tabular}
\end{table}

MORPH\_ELLIPSE at $(2\times2)$ was selected. The elliptical shape is preferable 
for natural images because it has no preferred axis direction, avoiding the 
directional bias of a rectangular kernel at larger sizes. At $2\times2$ the 
practical difference is negligible, but ELLIPSE was retained as best practice.

\section{Final Approach}
The final pipeline:

\begin{enumerate}
    \item \textbf{Light Gaussian pre-smoothing} $(3\times3)$, $\sigma=0.5$ — 
    attenuates the most extreme noise pixel amplitudes to improve NLM patch matching.
    \item \textbf{NLM denoising} (\texttt{h=10, hColor=10, templateWindowSize=7, 
    searchWindowSize=21}) — primary denoising stage, exploits non-local patch 
    self-similarity to remove dense noise while preserving colour gradients.
    \item \textbf{Morphological opening} (ELLIPSE $2\times2$) — removes residual 
    isolated noise clusters not fully handled by NLM.
\end{enumerate}

\begin{lstlisting}
import cv2 as cv

img = cv.imread("noisy.jpg")

# Stage 1: Light Gaussian pre-smoothing
gaus = cv.GaussianBlur(img, (3, 3), 0.5)

# Stage 2: Non-Local Means - primary denoising
result = cv.fastNlMeansDenoisingColored(
    gaus, None,
    h=10, hColor=10,
    templateWindowSize=7,
    searchWindowSize=21
)

# Stage 3: Morphological opening - residual noise cleanup
kernel = cv.getStructuringElement(cv.MORPH_ELLIPSE, (2, 2))
morph = cv.morphologyEx(result, cv.MORPH_OPEN, kernel)

cv.imwrite("scenary_final.png", morph)
\end{lstlisting}

\section{Results and Observations}
\begin{itemize}
    \item The sky gradient (cyan-to-blue) is smoothly restored with no banding 
    or posterisation artefacts.
    \item The teal horizontal band across the middle of the scene is clearly 
    visible and correctly coloured.
    \item The dark hill silhouette at the bottom is intact with sharp boundary 
    preservation.
    \item Residual very fine texture noise may be visible under high magnification 
    but is imperceptible at normal viewing scale.
    \item NLM's LAB-space operation is the critical factor for colour fidelity — 
    processing luminance and colour channels with appropriate separate strengths 
    prevents the hue shifts seen with naive per-channel median filtering.
\end{itemize}



\section{Conclusion}
For the colour scenery image with dense multicolour noise, Non-Local Means 
denoising with light Gaussian pre-smoothing and morphological post-processing 
was found to be the optimal pipeline. The key advantage of NLM over local 
filters is its large search window: in very high noise density conditions, 
local filters fail because the majority of nearby pixels are corrupted, 
whereas NLM finds clean matching patches from across the image. Parameter 
tuning revealed that \texttt{h=10, hColor=10, templateWindowSize=7, searchWindowSize=21} 
best balanced noise suppression against colour gradient and texture preservation.

\begin{thebibliography}{1}
\bibitem{opencv_freecodecamp}
Beau Carnes, https://www.freecodecamp.org/news/opencv-full-course/

\bibitem{opencv_docs}
https://docs.opencv.org/

\bibitem{geeksforgeeks}
https://www.geeksforgeeks.org/

\bibitem{GitHubCommands}
Adam Lusted (etoxin), https://gist.github.com/etoxin/1acb257550b1de60fe99

\end{thebibliography}

\end{document}